\section{Introducción}

El transporte cooperativo de cargas mediante múltiples robots móviles constituye un problema relevante en robótica distribuida por su relación con coordinación, estabilidad, reparto de esfuerzos y tolerancia a incertidumbres. En escenarios con varios AGVs, la necesidad de mantener coherencia cinemática y dinámica bajo información local hace especialmente pertinente el uso de estrategias distribuidas basadas en grafos y consenso \citep{bullo2009distributed,olfati2004consensus,ren2008distributed}.

Este TFM se centra en un escenario bidimensional con entre dos y cuatro AGVs que cooperan para desplazar una carga común. La motivación principal es estudiar cómo un controlador distribuido nominal puede extenderse hacia una versión adaptativa o robusta cuando la inercia efectiva de la carga no es conocida con precisión. El trabajo se orienta a simulación numérica reproducible, con interés en métricas cuantitativas de seguimiento, error de formación, esfuerzo de control, tiempo de convergencia y robustez ante perturbaciones.
